% ============================================================
%  BANK SOAL MATEMATIKA SMA — KURIKULUM MERDEKA
%  BAGIAN 1: FASE E  —  VERSI SOAL (REVISI)
%  - Tanpa banner; opsi A-E tersusun ke bawah
%  - Sinkron dengan Versi Lengkap (soal & nomor identik)
%  Kelas dokumen: exam  |  Kompilasi: pdfLaTeX (Overleaf)
% ============================================================
\documentclass[11pt,a4paper]{exam}

\usepackage[utf8]{inputenc}
\usepackage[T1]{fontenc}
\usepackage[bahasa]{babel}
\usepackage{amsmath,amssymb}
\usepackage{geometry}
\usepackage{xcolor}
\usepackage{tikz}
\usepackage{pgfplots}
\pgfplotsset{compat=1.18}
\usepackage{enumitem}
\usepackage{array,booktabs}

\geometry{a4paper,margin=2.3cm}

% (Versi Soal: TANPA \printanswers)

% ---------- WARNA ----------
\definecolor{biru}{HTML}{1F4E79}
\definecolor{hijau}{HTML}{2E7D32}
\definecolor{oranye}{HTML}{C75B12}
\definecolor{ungu}{HTML}{6A1B9A}

% ---------- HEADER / FOOTER ----------
\pagestyle{headandfoot}
\runningheadrule
\firstpageheadrule
\header{\textbf{\textcolor{biru}{BANK SOAL MATEMATIKA}}}{}{\textbf{FASE E}}
\footer{}{Halaman \thepage}{}

% ---------- LABEL PILIHAN: A. B. C. ----------
\renewcommand{\choicelabel}{\Alph{choice}.}

% ---------- JARAK ANTAR PILIHAN (tersusun ke bawah) ----------
\renewcommand{\choiceshook}{%
  \setlength{\leftmargin}{2em}%
  \setlength{\itemsep}{5pt}%
  \setlength{\topsep}{5pt}%
  \setlength{\parsep}{0pt}}

% ---------- TANDA TINGKAT KESULITAN ----------
\newcommand{\Fund}{{\footnotesize\textsf{\textcolor{hijau}{[Fundamental]}}}\ }
\newcommand{\Inter}{{\footnotesize\textsf{\textcolor{biru}{[Intermediate]}}}\ }
\newcommand{\Adv}{{\footnotesize\textsf{\textcolor{oranye}{[Advanced]}}}\ }
\newcommand{\Exp}{{\footnotesize\textsf{\textcolor{ungu}{[Expert]}}}\ }

% ---------- JUDUL TOPIK (heading sederhana, tanpa banner) ----------
\newcommand{\topik}[1]{%
  \par\addvspace{14pt}%
  \noindent{\large\bfseries\color{biru}#1}\par\addvspace{6pt}}

\begin{document}

\begin{center}
{\color{biru}\rule{\linewidth}{1.5pt}}\\[4pt]
{\LARGE\bfseries\color{biru} BANK SOAL MATEMATIKA SMA}\\[2pt]
{\large Versi Soal — Bagian 1: Fase E}\\[2pt]
{\small 65 Soal Pilihan Ganda}\\[4pt]
{\color{biru}\rule{\linewidth}{1.5pt}}
\end{center}
\vspace{4pt}

\noindent\textit{Petunjuk:} Pilih satu jawaban yang paling tepat. Label tingkat
kesulitan tertera pada awal tiap soal.

\topik{1. Eksponen dan Logaritma (12 Soal)}

\begin{questions}

\question \Fund Nilai dari $\dfrac{3^{7}}{3^{4}}$ adalah \dots
\begin{choices}
\choice $9$ \CorrectChoice $27$ \choice $81$ \choice $3$ \choice $343$
\end{choices}

\question \Fund Nilai dari $16^{3/4}$ adalah \dots
\begin{choices}
\CorrectChoice $8$ \choice $12$ \choice $6$ \choice $4$ \choice $64$
\end{choices}

\question \Fund Nilai dari ${}^{2}\!\log 32$ adalah \dots
\begin{choices}
\choice $4$ \CorrectChoice $5$ \choice $6$ \choice $8$ \choice $16$
\end{choices}

\question \Fund Nilai dari ${}^{3}\!\log 81$ adalah \dots
\begin{choices}
\choice $3$ \CorrectChoice $4$ \choice $5$ \choice $9$ \choice $27$
\end{choices}

\question \Inter Tabel berikut menunjukkan populasi bakteri tiap jam.

\begin{center}\small
\begin{tabular}{|c|c|c|c|c|}
\hline
Jam ($t$) & 0 & 1 & 2 & 3\\\hline
Populasi & 100 & 300 & 900 & 2700\\\hline
\end{tabular}
\end{center}

Jika pola berlanjut, populasi pada jam ke-5 adalah \dots
\begin{choices}
\choice $8100$ \choice $13500$ \CorrectChoice $24300$ \choice $72900$ \choice $2700$
\end{choices}

\question \Inter Nilai dari ${}^{2}\!\log 6 + {}^{2}\!\log 4 - {}^{2}\!\log 3$ adalah \dots
\begin{choices}
\choice $2$ \CorrectChoice $3$ \choice $4$ \choice $5$ \choice $1$
\end{choices}

\question \Inter Bentuk sederhana dari $\dfrac{(x^{2}y^{-3})^{2}}{x^{-1}y}$ adalah \dots
\begin{choices}
\choice $\dfrac{x^{3}}{y^{5}}$ \CorrectChoice $\dfrac{x^{5}}{y^{7}}$
\choice $x^{5}y^{7}$ \choice $\dfrac{x^{4}}{y^{6}}$ \choice $\dfrac{x^{5}}{y^{6}}$
\end{choices}

\question \Inter Diketahui ${}^{a}\!\log 2 = 0{,}3$ dan ${}^{a}\!\log 3 = 0{,}5$.
Nilai ${}^{a}\!\log 12$ adalah \dots
\begin{choices}
\choice $0{,}8$ \choice $0{,}9$ \choice $1{,}0$ \CorrectChoice $1{,}1$ \choice $1{,}3$
\end{choices}

\question \Adv Jumlah semua nilai $x$ yang memenuhi $4^{x} - 5\cdot 2^{x} + 4 = 0$ adalah \dots
\begin{choices}
\choice $1$ \CorrectChoice $2$ \choice $3$ \choice $4$ \choice $5$
\end{choices}

\question \Adv Grafik fungsi $f(x) = a^{x}$ ($a>1$) melalui titik $(2,9)$.

\begin{center}
\begin{tikzpicture}
\begin{axis}[axis lines=middle,width=6.5cm,height=4.2cm,
  xmin=-1,xmax=3,ymin=-1,ymax=10,xtick={2},ytick={9},
  xticklabel style={font=\small},yticklabel style={font=\small}]
\addplot[very thick,biru,domain=-1:2.1,samples=40]{3^x};
\addplot[only marks,mark=*,oranye,mark size=2pt]coordinates{(2,9)};
\end{axis}
\end{tikzpicture}
\end{center}

Nilai $f(4)$ adalah \dots
\begin{choices}
\choice $27$ \CorrectChoice $81$ \choice $243$ \choice $9$ \choice $12$
\end{choices}

\question \Adv Populasi kota dimodelkan $P(t) = 5000\,(1{,}2)^{t}$ ($t$ tahun).
Tahun pertama populasi melampaui 8000 adalah pada \dots
\begin{choices}
\choice tahun ke-2 \CorrectChoice tahun ke-3 \choice tahun ke-4
\choice tahun ke-5 \choice tahun ke-6
\end{choices}

\question \Exp Jika ${}^{2}\!\log 3 = p$, maka ${}^{6}\!\log 8$ dinyatakan dalam $p$ adalah \dots
\begin{choices}
\CorrectChoice $\dfrac{3}{1+p}$ \choice $\dfrac{3}{p+2}$ \choice $3p$
\choice $\dfrac{p}{3}$ \choice $\dfrac{1+p}{3}$
\end{choices}

\topik{2. Barisan dan Deret (10 Soal)}

\question \Fund Suku ke-10 dari barisan aritmetika $4, 7, 10, \dots$ adalah \dots
\begin{choices}
\choice $28$ \choice $30$ \CorrectChoice $31$ \choice $33$ \choice $34$
\end{choices}

\question \Fund Rasio dari barisan geometri $2, 6, 18, \dots$ adalah \dots
\begin{choices}
\choice $2$ \CorrectChoice $3$ \choice $4$ \choice $6$ \choice $\tfrac{1}{3}$
\end{choices}

\question \Fund Jumlah $3 + 6 + 9 + 12 + 15$ adalah \dots
\begin{choices}
\choice $30$ \choice $40$ \CorrectChoice $45$ \choice $50$ \choice $75$
\end{choices}

\question \Inter Pada barisan aritmetika diketahui $U_3 = 10$ dan $U_7 = 26$.
Beda barisan tersebut adalah \dots
\begin{choices}
\choice $2$ \choice $3$ \CorrectChoice $4$ \choice $5$ \choice $6$
\end{choices}

\question \Inter Jumlah deret geometri tak hingga $9 + 3 + 1 + \cdots$ adalah \dots
\begin{choices}
\choice $12$ \choice $13$ \CorrectChoice $13{,}5$ \choice $18$ \choice $27$
\end{choices}

\question \Inter Seorang siswa menabung Rp2.000 pada hari pertama dan menambah
tabungan Rp2.000 lebih banyak setiap hari berikutnya. Total tabungan selama
10 hari adalah \dots
\begin{choices}
\choice Rp90.000 \CorrectChoice Rp110.000 \choice Rp100.000
\choice Rp120.000 \choice Rp55.000
\end{choices}

\question \Adv Sebuah tali dipotong menjadi 5 bagian yang membentuk barisan
geometri. Jika potongan terpendek 6 cm dan terpanjang 96 cm, panjang tali
semula adalah \dots
\begin{choices}
\choice $150$ cm \choice $168$ cm \choice $180$ cm \CorrectChoice $186$ cm \choice $192$ cm
\end{choices}

\question \Adv Sebuah bola dijatuhkan dari ketinggian 10 m dan setiap pantulan
mencapai $\tfrac{3}{4}$ tinggi sebelumnya. Panjang lintasan total hingga bola
berhenti adalah \dots
\begin{choices}
\choice $40$ m \choice $60$ m \CorrectChoice $70$ m \choice $80$ m \choice $100$ m
\end{choices}

\question \Adv Jumlah $n$ suku pertama suatu deret adalah $S_n = 2n^{2} + 3n$.
Suku ke-5 deret tersebut adalah \dots
\begin{choices}
\choice $17$ \choice $20$ \CorrectChoice $21$ \choice $23$ \choice $25$
\end{choices}

\question \Exp Jumlah tak hingga suatu deret geometri adalah 12 dan suku
pertamanya 8. Suku ke-3 deret tersebut adalah \dots
\begin{choices}
\choice $\tfrac{8}{3}$ \choice $\tfrac{4}{3}$ \CorrectChoice $\tfrac{8}{9}$
\choice $\tfrac{2}{3}$ \choice $24$
\end{choices}

\topik{3. Vektor (8 Soal)}

\question \Fund Panjang vektor $\vec{a} = (3,4)$ adalah \dots
\begin{choices}
\CorrectChoice $5$ \choice $7$ \choice $12$ \choice $25$ \choice $\sqrt{7}$
\end{choices}

\question \Fund Diketahui $\vec{a} = (2,1)$ dan $\vec{b} = (1,3)$. Hasil $\vec{a} + \vec{b}$ adalah \dots
\begin{choices}
\choice $(1,4)$ \choice $(2,3)$ \CorrectChoice $(3,4)$ \choice $(3,3)$ \choice $(1,-2)$
\end{choices}

\question \Inter Diketahui $\vec{a} = (2,-1)$ dan $\vec{b} = (3,4)$. Hasil $\vec{a}\cdot\vec{b}$ adalah \dots
\begin{choices}
\choice $-2$ \CorrectChoice $2$ \choice $6$ \choice $10$ \choice $14$
\end{choices}

\question \Inter Diketahui $\vec{a} = (1,2)$ dan $\vec{b} = (3,-1)$. Hasil $2\vec{a} - \vec{b}$ adalah \dots
\begin{choices}
\CorrectChoice $(-1,5)$ \choice $(5,-1)$ \choice $(-1,3)$ \choice $(1,5)$ \choice $(-5,1)$
\end{choices}

\question \Inter Vektor $\vec{a} = (2,k)$ tegak lurus terhadap $\vec{b} = (3,-6)$. Nilai $k$ adalah \dots
\begin{choices}
\CorrectChoice $1$ \choice $-1$ \choice $2$ \choice $3$ \choice $4$
\end{choices}

\question \Adv Besar sudut antara vektor $\vec{a} = (1,0)$ dan $\vec{b} = (1,1)$ adalah \dots
\begin{choices}
\choice $0^\circ$ \choice $30^\circ$ \CorrectChoice $45^\circ$ \choice $60^\circ$ \choice $90^\circ$
\end{choices}

\question \Adv Dua gaya $\vec{F_1}$ dan $\vec{F_2}$ bekerja pada satu titik dengan
$|\vec{F_1}|=3$ N, $|\vec{F_2}|=4$ N, dan sudut antara keduanya $60^\circ$.
Besar resultannya adalah \dots
\begin{choices}
\choice $5$ N \choice $7$ N \choice $\sqrt{13}$ N \CorrectChoice $\sqrt{37}$ N \choice $\sqrt{61}$ N
\end{choices}

\question \Exp Diketahui titik $A(0,0)$, $B(2,0)$, dan $C(1,\sqrt{3})$. Besar
sudut $BAC$ adalah \dots
\begin{choices}
\choice $30^\circ$ \CorrectChoice $60^\circ$ \choice $45^\circ$
\choice $90^\circ$ \choice $120^\circ$
\end{choices}

\topik{4. Trigonometri Dasar (10 Soal)}

\question \Fund Nilai $\sin 30^\circ$ adalah \dots
\begin{choices}
\CorrectChoice $\tfrac{1}{2}$ \choice $\tfrac{\sqrt{2}}{2}$ \choice $\tfrac{\sqrt{3}}{2}$ \choice $1$ \choice $0$
\end{choices}

\question \Fund Nilai $\cos 60^\circ$ adalah \dots
\begin{choices}
\CorrectChoice $\tfrac{1}{2}$ \choice $\tfrac{\sqrt{2}}{2}$ \choice $\tfrac{\sqrt{3}}{2}$ \choice $1$ \choice $0$
\end{choices}

\question \Fund Nilai $\tan 45^\circ$ adalah \dots
\begin{choices}
\choice $0$ \CorrectChoice $1$ \choice $\sqrt{3}$ \choice $\tfrac{\sqrt{3}}{3}$ \choice $\tfrac{1}{2}$
\end{choices}

\question \Fund Pada segitiga siku-siku, sisi depan sudut $\alpha$ adalah 3 dan sisi
miring 5. Nilai $\sin\alpha$ adalah \dots
\begin{choices}
\CorrectChoice $\tfrac{3}{5}$ \choice $\tfrac{4}{5}$ \choice $\tfrac{3}{4}$ \choice $\tfrac{5}{3}$ \choice $\tfrac{4}{3}$
\end{choices}

\question \Inter Pada segitiga siku-siku, sisi samping $= 8$ dan sisi miring $= 10$.
Nilai $\cos\alpha$ adalah \dots
\begin{choices}
\choice $\tfrac{3}{5}$ \CorrectChoice $\tfrac{4}{5}$ \choice $\tfrac{3}{4}$ \choice $\tfrac{5}{4}$ \choice $\tfrac{4}{3}$
\end{choices}

\question \Inter Diketahui $\sin\alpha = \tfrac{3}{5}$ dengan $\alpha$ sudut lancip.
Nilai $\cos\alpha$ adalah \dots
\begin{choices}
\CorrectChoice $\tfrac{4}{5}$ \choice $\tfrac{3}{4}$ \choice $\tfrac{5}{4}$ \choice $\tfrac{3}{5}$ \choice $\tfrac{4}{3}$
\end{choices}

\question \Inter Sebuah tangga sepanjang 5 m bersandar pada dinding membentuk
sudut $60^\circ$ dengan tanah. Tinggi dinding yang dicapai ujung tangga adalah \dots
\begin{choices}
\CorrectChoice $\tfrac{5\sqrt{3}}{2}$ m \choice $\tfrac{5}{2}$ m
\choice $\tfrac{5\sqrt{2}}{2}$ m \choice $5$ m \choice $5\sqrt{3}$ m
\end{choices}

\question \Adv Dari sebuah titik yang berjarak 20 m dari kaki gedung, sudut elevasi
puncak gedung adalah $30^\circ$. Tinggi gedung tersebut adalah \dots
\begin{choices}
\choice $10$ m \CorrectChoice $\tfrac{20\sqrt{3}}{3}$ m \choice $10\sqrt{3}$ m
\choice $20\sqrt{3}$ m \choice $20$ m
\end{choices}

\question \Adv Nilai dari $\sin 30^\circ \cos 60^\circ + \cos 30^\circ \sin 60^\circ$ adalah \dots
\begin{choices}
\choice $0$ \choice $\tfrac{1}{2}$ \choice $\tfrac{\sqrt{3}}{2}$ \CorrectChoice $1$ \choice $\sqrt{3}$
\end{choices}

\question \Exp Dua pengamat berdiri terpisah 100 m. Sebuah balon di antara keduanya
terlihat dengan sudut elevasi $30^\circ$ (dari pengamat pertama) dan $45^\circ$
(dari pengamat kedua). Tinggi balon adalah \dots
\begin{choices}
\choice $25$ m \choice $50$ m \CorrectChoice $50(\sqrt{3}-1)$ m
\choice $50(\sqrt{3}+1)$ m \choice $\tfrac{100}{\sqrt{3}}$ m
\end{choices}

\topik{5. Sistem Persamaan \& Pertidaksamaan Linear (8 Soal)}

\question \Fund Penyelesaian sistem $x + y = 5$ dan $x - y = 1$ memberikan nilai $x = \dots$
\begin{choices}
\choice $1$ \choice $2$ \CorrectChoice $3$ \choice $4$ \choice $5$
\end{choices}

\question \Fund Diketahui $2x + y = 7$ dan $x = 2$. Nilai $y$ adalah \dots
\begin{choices}
\choice $1$ \choice $2$ \CorrectChoice $3$ \choice $4$ \choice $5$
\end{choices}

\question \Inter Penyelesaian SPLTV $\begin{cases} x+y+z=6\\ x-y+z=2\\ x+y-z=0 \end{cases}$
memberikan nilai $z = \dots$
\begin{choices}
\choice $1$ \choice $2$ \CorrectChoice $3$ \choice $4$ \choice $5$
\end{choices}

\question \Inter Penyelesaian $x + 2y = 8$ dan $3x - y = 3$ memberikan nilai $x + y = \dots$
\begin{choices}
\choice $3$ \choice $4$ \CorrectChoice $5$ \choice $7$ \choice $2$
\end{choices}

\question \Inter Perhatikan daerah penyelesaian (DP) yang diarsir berikut, dibatasi
$x+y\le 4$, $x\ge 0$, $y\ge 0$.

\begin{center}
\begin{tikzpicture}[scale=0.65]
\fill[biru!12] (0,0)--(4,0)--(0,4)--cycle;
\draw[->] (-0.5,0)--(5,0) node[right]{$x$};
\draw[->] (0,-0.5)--(0,5) node[above]{$y$};
\draw[very thick,biru] (4,0)--(0,4) node[above right]{\small $x+y=4$};
\node at (1,1) {\small DP};
\end{tikzpicture}
\end{center}

Titik berikut yang berada di dalam daerah penyelesaian adalah \dots
\begin{choices}
\CorrectChoice $(1,1)$ \choice $(5,0)$ \choice $(3,3)$ \choice $(0,5)$ \choice $(2,3)$
\end{choices}

\question \Adv Harga 2 buku dan 3 pensil Rp13.000, sedangkan 3 buku dan 2 pensil
Rp12.000. Harga satu buku adalah \dots
\begin{choices}
\choice Rp1.500 \CorrectChoice Rp2.000 \choice Rp2.500 \choice Rp3.000 \choice Rp4.000
\end{choices}

\question \Adv Tiga bilangan berjumlah 33. Bilangan terbesar dua kali bilangan
terkecil, dan bilangan tengah lima lebihnya dari bilangan terkecil. Bilangan
terbesar adalah \dots
\begin{choices}
\choice $7$ \choice $12$ \CorrectChoice $14$ \choice $15$ \choice $16$
\end{choices}

\question \Exp Sistem $x + y = 4$ dan $ax + y = 6$ tidak memiliki penyelesaian jika
nilai $a$ sama dengan \dots
\begin{choices}
\choice $-1$ \choice $0$ \CorrectChoice $1$ \choice $2$ \choice $4$
\end{choices}

\topik{6. Fungsi Kuadrat (8 Soal)}

\question \Fund Jumlah akar-akar persamaan $x^{2} - 5x + 6 = 0$ adalah \dots
\begin{choices}
\choice $-5$ \choice $1$ \CorrectChoice $5$ \choice $6$ \choice $-1$
\end{choices}

\question \Fund Hasil kali akar-akar persamaan $x^{2} - 5x + 6 = 0$ adalah \dots
\begin{choices}
\choice $-6$ \choice $-5$ \choice $5$ \CorrectChoice $6$ \choice $1$
\end{choices}

\question \Fund Sumbu simetri grafik $f(x) = x^{2} - 4x + 1$ adalah garis \dots
\begin{choices}
\choice $x = -2$ \choice $x = -4$ \choice $x = 1$ \CorrectChoice $x = 2$ \choice $x = 4$
\end{choices}

\question \Inter Titik puncak grafik $f(x) = x^{2} - 4x + 3$ adalah \dots
\begin{choices}
\CorrectChoice $(2,-1)$ \choice $(2,1)$ \choice $(-2,-1)$ \choice $(2,-3)$ \choice $(1,2)$
\end{choices}

\question \Inter Berdasarkan diskriminan, persamaan $2x^{2} + 3x + 5 = 0$ memiliki \dots
\begin{choices}
\choice dua akar real berbeda
\CorrectChoice tidak memiliki akar real (akar imajiner)
\choice satu akar real (kembar)
\choice dua akar real positif
\choice dua akar real negatif
\end{choices}

\question \Inter Sebuah parabola terbuka ke bawah, berpuncak di $(1,4)$, dan melalui
titik $(0,3)$. Persamaannya adalah \dots
\begin{choices}
\CorrectChoice $f(x) = -x^{2} + 2x + 3$
\choice $f(x) = x^{2} - 2x + 3$
\choice $f(x) = -x^{2} - 2x + 3$
\choice $f(x) = -x^{2} + 2x - 3$
\choice $f(x) = x^{2} + 2x + 3$
\end{choices}

\question \Adv Persamaan $x^{2} + mx + 9 = 0$ memiliki akar kembar untuk nilai $m$
positif sebesar \dots
\begin{choices}
\choice $3$ \CorrectChoice $6$ \choice $9$ \choice $12$ \choice $18$
\end{choices}

\question \Exp Jika $p$ dan $q$ adalah akar-akar persamaan $x^{2} - 3x + 1 = 0$, maka
nilai $p^{2} + q^{2}$ adalah \dots
\begin{choices}
\choice $3$ \choice $5$ \CorrectChoice $7$ \choice $9$ \choice $11$
\end{choices}

\topik{7. Statistika (4 Soal)}

\question \Fund Rata-rata (mean) dari data $5, 7, 9, 11$ adalah \dots
\begin{choices}
\choice $6$ \choice $7$ \CorrectChoice $8$ \choice $9$ \choice $32$
\end{choices}

\question \Inter Tabel berikut menunjukkan nilai ulangan 10 siswa.

\begin{center}\small
\begin{tabular}{|c|c|c|c|c|}
\hline
Nilai & 60 & 70 & 80 & 90\\\hline
Frekuensi & 2 & 3 & 4 & 1\\\hline
\end{tabular}
\end{center}

Rata-rata nilai tersebut adalah \dots
\begin{choices}
\choice $70$ \CorrectChoice $74$ \choice $75$ \choice $72$ \choice $80$
\end{choices}

\question \Inter Modus dari data $4, 6, 7, 7, 9, 7, 4$ adalah \dots
\begin{choices}
\choice $4$ \choice $6$ \CorrectChoice $7$ \choice $9$ \choice $5$
\end{choices}

\question \Adv Rata-rata 5 nilai ulangan adalah 70. Setelah ditambah satu nilai,
rata-rata 6 nilai menjadi 72. Nilai yang baru ditambahkan adalah \dots
\begin{choices}
\choice $72$ \choice $75$ \choice $80$ \CorrectChoice $82$ \choice $90$
\end{choices}

\topik{8. Peluang (5 Soal)}

\question \Fund Sebuah dadu dilempar sekali. Peluang muncul mata dadu 5 adalah \dots
\begin{choices}
\CorrectChoice $\tfrac{1}{6}$ \choice $\tfrac{1}{3}$ \choice $\tfrac{1}{2}$ \choice $\tfrac{5}{6}$ \choice $1$
\end{choices}

\question \Fund Sebuah koin dilempar sekali. Peluang muncul sisi angka adalah \dots
\begin{choices}
\choice $\tfrac{1}{4}$ \choice $\tfrac{1}{3}$ \CorrectChoice $\tfrac{1}{2}$ \choice $1$ \choice $2$
\end{choices}

\question \Inter Sebuah dadu dilempar. Peluang muncul mata dadu genap atau lebih
dari 4 adalah \dots
\begin{choices}
\choice $\tfrac{1}{3}$ \choice $\tfrac{1}{2}$ \CorrectChoice $\tfrac{2}{3}$ \choice $\tfrac{5}{6}$ \choice $\tfrac{3}{6}$
\end{choices}

\question \Inter Sebuah kantong berisi 5 bola merah dan 3 bola putih. Satu bola
diambil secara acak. Peluang terambil bola putih adalah \dots
\begin{choices}
\CorrectChoice $\tfrac{3}{8}$ \choice $\tfrac{5}{8}$ \choice $\tfrac{3}{5}$ \choice $\tfrac{1}{3}$ \choice $\tfrac{8}{3}$
\end{choices}

\question \Adv Dua buah dadu dilempar bersamaan. Peluang jumlah kedua mata dadu
sama dengan 7 adalah \dots
\begin{choices}
\CorrectChoice $\tfrac{1}{6}$ \choice $\tfrac{1}{9}$ \choice $\tfrac{5}{36}$ \choice $\tfrac{7}{36}$ \choice $\tfrac{1}{12}$
\end{choices}

\end{questions}

\vspace{6pt}
\begin{center}
{\color{biru}\rule{\linewidth}{1pt}}\\[3pt]
{\small\textit{Akhir Bagian 1 (Fase E) — Versi Soal. Kunci jawaban \& pembahasan ada pada Versi Lengkap.}}
\end{center}

\end{document}